% GENERATED FILE. Edit canonical lessons, then run npm run generate:books.
% canonical-source-hash: fnv1a64:613c242abb4fddb6
% canonical-lessons: ML-C46-fruit, ML-C46-cloth, ML-C46-lamp, ML-C46-salt, ML-C46-book

\chapter{Things You Ask For}
\label{ch:ml-ask-for}
\hlchaptermodality{46}

\begin{chapteropening}
\textbf{By the end of this chapter:} \emph{I can name five everyday things I might ask someone for in Malayalam, and say which of them the language has held all along and which walked in from elsewhere.}

Complete ML-C46-book and say all five words in order.
\end{chapteropening}

\section[paḻaṁ]{\ml{പഴം} (paḻaṁ) — a fruit}

\label{lesson:ML-C46-fruit}



\begin{quote}
Before the new run of words: what did \emph{bharttāvŭ} mean, and what did \emph{bhārya} mean?
\end{quote}

\subsection*{You'll want to know: \ml{പഴം}}
\textbf{\ml{പഴം}} (\emph{paḻaṁ}) — \textquotedblleft{}a fruit\textquotedblright{}.

Behind the fruit sits a verb of ripening. Proto-Dravidian \textbf{*paḻ}, 'to ripen, to grow old', handed this family both a fruit and an age, and Malayalam kept the pair close enough that a learner can still hear it: \ml{പഴയ} (\emph{paḻaya}), the \textquotedblleft{}old\textquotedblright{} from the run of adjectives, is that same root wearing an adjective ending.

Tamil says \ta{பழம்} (\emph{paḻam}), which is very nearly the identical word — the ordinary state of affairs between these two closest sisters. Telugu went to \te{పండు} (\emph{paṇḍu}) and Kannada to \ka{ಹಣ್ಣು} (\emph{haṇṇu}), and the Kannada form has lost its opening \emph{p-} to a shift Malayalam sat out entirely.

The \ml{ഴ} in the middle is the letter this book has already called Malayalam's own oddity, the one only these two sisters still write.

\begin{grammarlens}[title={What you've built}]
The first of five things you might ask someone for.
\end{grammarlens}

\subsection*{Your turn}
Say these aloud:

\begin{itemize}
\raggedright
  \item \emph{paḻaṁ}
  \item it once more, slowly
  \item \emph{paḻaṁ}, then \emph{dayavāyi paḻaṁ tarū} — \textquotedblleft{}please give a fruit\textquotedblright{}
\end{itemize}

\begin{tcolorbox}[breakable,colback=teal!4,colframe=teal!35!black,title={Before you move on}]
What does \ml{പഴം} mean? (\textquotedblleft{}a fruit\textquotedblright{}.) Say it once more.
\end{tcolorbox}
\section[tuṇi]{\ml{തുണി} (tuṇi) — cloth}

\label{lesson:ML-C46-cloth}



\begin{quote}
Before the new one: what did \emph{paḻaṁ} mean?
\end{quote}

\subsection*{You'll want to know: \ml{തുണി}}
\textbf{\ml{തുണി}} (\emph{tuṇi}) — \textquotedblleft{}cloth\textquotedblright{}.

\ml{തുണി} is the plain, unceremonious word: the bolt on the shelf, the washing on the line, the rag in your hand. Tamil writes \ta{துணி} (\emph{tuṇi}) and means the same thing.

Malayalam holds a dressier word for the same stuff — \ml{വസ്ത്രം} (\emph{vastraṁ}), out of Sanskrit — and it turns up on a wedding invitation rather than over a washing line. The native word does the daily work; the borrowed one does the ceremony.

That division of labour is worth watching for, because it repeats through this book: two words, one meaning, and the register telling them apart.

\begin{grammarlens}[title={What you've built}]
Two: a fruit and a cloth.
\end{grammarlens}

\subsection*{Your turn}
Say these aloud:

\begin{itemize}
\raggedright
  \item \emph{tuṇi}
  \item it once more, slowly
  \item \emph{paḻaṁ}, then \emph{tuṇi}, and ask for each one with \emph{dayavāyi}
\end{itemize}

\begin{tcolorbox}[breakable,colback=teal!4,colframe=teal!35!black,title={Before you move on}]
What does \ml{തുണി} mean? (\textquotedblleft{}cloth\textquotedblright{}.) Say it once more.
\end{tcolorbox}
\section[viḷakkŭ]{\ml{വിളക്ക്} (viḷakkŭ) — a lamp}

\label{lesson:ML-C46-lamp}



\begin{quote}
Before the new one: what did \emph{tuṇi} mean?
\end{quote}

\subsection*{You'll want to know: \ml{വിളക്ക്}}
\textbf{\ml{വിളക്ക്}} (\emph{viḷakkŭ}) — \textquotedblleft{}a lamp\textquotedblright{}.

A native word, built on a root about shining. Tamil \ta{விளக்கு} (\emph{viḷakku}) is the same word, and in both languages a related sense — making something clear, shining a light on it — sits alongside the object.

The object itself has a particular shape in Kerala: the \ml{നിലവിളക്ക്} (\emph{nilaviḷakkŭ}), the tall standing brass lamp lit at a doorway on an occasion. Its first piece is \ml{നില} (\emph{nila}), the ground it stands on.

Kannada reached for a borrowed word here, \ka{ದೀಪ} (\emph{dīpa}) from Sanskrit, where Malayalam kept its own. Malayalam has \ml{ദീപം} (\emph{dīpaṁ}) as well, in temple and festival senses, which leaves the everyday lamp to the older resident.

\begin{grammarlens}[title={What you've built}]
Three. A fruit, a cloth, a lamp.
\end{grammarlens}

\subsection*{Your turn}
Say these aloud:

\begin{itemize}
\raggedright
  \item \emph{viḷakkŭ}
  \item it once more, slowly
  \item \emph{viḷakkŭ}, then \emph{tuṇi}, and hear which of the two is the older resident
\end{itemize}

\begin{tcolorbox}[breakable,colback=teal!4,colframe=teal!35!black,title={Before you move on}]
What does \ml{വിളക്ക്} mean? (\textquotedblleft{}a lamp\textquotedblright{}.) Say it once more.
\end{tcolorbox}
\section[uppŭ]{\ml{ഉപ്പ്} (uppŭ) — salt}

\label{lesson:ML-C46-salt}



\begin{quote}
Before the new one: what did \emph{viḷakkŭ} mean?
\end{quote}

\subsection*{You'll want to know: \ml{ഉപ്പ്}}
\textbf{\ml{ഉപ്പ്}} (\emph{uppŭ}) — \textquotedblleft{}salt\textquotedblright{}.

Proto-Dravidian \textbf{*uppu}, and the whole family kept it very nearly untouched: Tamil \ta{உப்பு}, Telugu \te{ఉప్పు}, Kannada \ka{ಉಪ್ಪು}. Four sisters, one word, four spellings of it.

There is a \emph{p} sitting inside this word, and in Kannada it did not turn into an \emph{h} the way the \emph{p} of its fruit-word did. The shift that separated \ka{ಹಣ್ಣು} from \ml{പഴം} reached only the FIRST sound of a word. Inside one, the \emph{p} stays where it was put — which is why salt looks so alike on all four sides of the family.

The final \ml{്} here is the half-vowel this book romanizes \emph{-ŭ}: a consonant let go of rather than cut off.

\begin{grammarlens}[title={What you've built}]
Four. Salt, and the fruit, and the cloth, and the lamp.
\end{grammarlens}

\subsection*{Your turn}
Say these aloud:

\begin{itemize}
\raggedright
  \item \emph{uppŭ}
  \item it once more, slowly
  \item \emph{uppŭ}, then \emph{paḻaṁ}, and notice the \emph{p} the cousins kept in one word and lost in the other
\end{itemize}

\begin{tcolorbox}[breakable,colback=teal!4,colframe=teal!35!black,title={Before you move on}]
What does \ml{ഉപ്പ്} mean? (\textquotedblleft{}salt\textquotedblright{}.) Say it once more.
\end{tcolorbox}
\section[pustakaṁ]{\ml{പുസ്തകം} (pustakaṁ) — a book}

\label{lesson:ML-C46-book}



\begin{quote}
Before the new one: what did \emph{uppŭ} mean?
\end{quote}

\subsection*{You'll want to know: \ml{പുസ്തകം}}
\textbf{\ml{പുസ്തകം}} (\emph{pustakaṁ}) — \textquotedblleft{}a book\textquotedblright{}.

Sanskrit \emph{pustaka}, taken whole. Tamil \ta{புத்தகம்}, Telugu \te{పుస్తకం} and Kannada \ka{ಪುಸ್ತಕ} are the same borrowing, so the four sisters agree here for a different reason than they agree about salt: not descent, but a shared debt.

Sanskrit itself seems to have taken it from further west. The usual trail runs to Persian \emph{post}, a hide or a skin — the material a book was, before it was paper. A word for the animal became a word for the page.

Malayalam also says \emph{granthaṁ} for a book of weight and age, another Sanskrit word, and that one is built on a root about knotting: palm leaves tied into a bundle. Two borrowings, two different pictures of what a book physically was.

\begin{grammarlens}[title={What you've built}]
Five, and the run is closed: a fruit, a cloth, a lamp, salt, a book.
\end{grammarlens}

\subsection*{Your turn}
Say these aloud:

\begin{itemize}
\raggedright
  \item \emph{pustakaṁ}
  \item it once more, slowly
  \item all five in order — \emph{paḻaṁ}, \emph{tuṇi}, \emph{viḷakkŭ}, \emph{uppŭ}, \emph{pustakaṁ}
\end{itemize}

\begin{tcolorbox}[breakable,colback=teal!4,colframe=teal!35!black,title={Before you move on}]
What does \ml{പുസ്തകം} mean? (\textquotedblleft{}a book\textquotedblright{}.) Say it once more.
\end{tcolorbox}
